% Example input for andre_style_check.py.
% The study is invented: a test-beam measurement of the time resolution of a small
% scintillator timing detector. Every number in it is made up.
% The text is written on purpose in a style that is not his: short sentences, "We" at
% the start of most of them, ", so" as the causal join, long captions, "see Fig.".

\section{Setup}
\label{sec:setup}

We test a small scintillator tile in a 120~GeV pion beam.
We read out the tile with two silicon photomultipliers.
We place the tile 2~m downstream of a reference counter.
We digitize both signals at 5~GS/s.
The trigger is a coincidence of two beam counters.
The beam spot is wider than the tile, so we select tracks that point to its central region.
For each run, we record approximately 50,000 events.

We utilize a constant-fraction threshold to extract the arrival time; see Fig.~\ref{fig:pulse}.
The threshold is 20\% of the pulse amplitude.
With this choice, the time walk is very small.
It should be noted that we do not correct for the residual walk.
We check the baseline before each pulse.
We reject events with a second pulse.
We align the tile with the beam axis by hand.

\begin{figure}
\includegraphics[width=\columnwidth]{pulse}
\caption{The pulse shape of the tile for 120~GeV pions. We fit the leading edge with a
straight line. The fit is very stable because the pulses are large. The red line shows
the fit and the blue points show the samples. As can be seen, the residual time walk is
small.}
\label{fig:pulse}
\end{figure}

\section{Time resolution}
\label{sec:resolution}

We average the times of the two sensors, so the result does not depend on the impact point.
We fit the time difference to the reference counter with a Gaussian.
We observe a width of approximately 40~ps.
The reference counter is much faster than the tile, so we neglect its jitter.
In addition, we repeat the fit at five bias voltages (Fig.~\ref{fig:voltage}).
Additionally, we leverage the track position to study the edges of the tile, as described
in Sec.~\ref{sec:setup}.

The resolution improves with the voltage.
It reaches 38~ps at 58~V.
We also vary the threshold between 10\% and 30\%.
We find no change in the resolution.
The light yield is lower near the edges, so the resolution is worse there.
As a result, we restrict the measurement to the central region.
We note that the resolution that we obtain after this restriction is significantly better
than the one that we obtained with the full tile and that it stays below the requirement
of the experiment over the whole range of bias voltages that we scanned in this test beam.

\begin{figure}
\includegraphics[width=\columnwidth]{voltage}
\caption{This figure shows the time resolution for different bias voltages. The resolution
improves with the voltage because the gain of the sensors rises. At the highest voltages
the curve flattens because the dark-count rate rises. The best value is 38~ps, which is an
excellent result.}
\label{fig:voltage}
\end{figure}

\section{Conclusions}

We measured the time resolution of a small scintillator tile in a test beam.
We find 38~ps at the best bias voltage.
The requirement is 50~ps, so the tile is suitable.
Importantly, the so-called constant-fraction method is robust.
Our state-of-the-art readout outperforms the previous one.
We see no dependence on the beam intensity.
We also tested a second tile.
We find the same resolution for it.
For the next test beam, we plan a larger tile, a thicker foil, etc.
